% Folder 136, batch 04 (archivists' pages 61-80).
% Pass: Opus 5.5 (claude-opus-5-5), 2026-10-03 — first pass, unchecked against the pages by a human.
%
% Pages 62, 70 and 72 are typescript pages of an unrelated text (verso), and page 78 is a
% blank folder cover: none is transcribed.
\documentclass[11pt,a4paper]{article}
\input{../preamble/grothendieck}

\folder{136}
\batch{4}
\pages{61}{80}
\foldertitle{Complexe de De Rham à puissance divisée [conférence de 1976 à l’IHÉS] : notes manuscrites (s.d.).}
\dating{[à partir de 1975-1976]}
\watermark{Édition de démonstration}

\begin{document}

\page{61}
\note{la moitié supérieure de la page est barrée de longs traits obliques ; on en donne ce qui se lit.}

\struck{$\Delta/\mathcal{X} \xrightarrow{\text{fibration}} \Delta$}

\struck{$C \to B$, \quad $F : C^{\circ} \to \mathcal{E}$, \quad $C_b \overset{i^*}{\leftarrow} C_{b'}$, \quad $x' \mapsto b'$ :}
\[
  \text{\struck{$(b \in B) \longmapsto \sum_{x \in C_b} F(x) \longrightarrow F(i_*(x'))$}}
\]
\[
  \text{\struck{$(b' \in B) \longmapsto \sum_{x' \in C_{b'}} F(x') \longrightarrow F(x')$}}
\]
\struck{(si dual)}
\[
  \text{\struck{$(b \in B) \longmapsto \prod_{x \in C_b} F(x) \longrightarrow F(i^*(x'))$}}
\]
\[
  \text{\struck{$(b' \in B) \longmapsto \prod_{x \in C_{b'}} F(x) \longrightarrow F(x')$}}
\]
\struck{$\mathcal{E}$ (cat. avec sommes directes)}

Si $(\Delta/\mathcal{X})^{\circ} \xrightarrow{A_*} \mathcal{E}$ \struck{\ill{}} est un faisceau sur
$\mathcal{X}$ (i.e. \struck{faisceau} \add{système} à coeff. contravariant par rapport au
simplexe variable sur $\mathcal{X}$), on lui associe $C_*(\mathcal{X}, A_*) = i_!(A_*) :
\Delta^{\circ} \to \mathcal{E}$, où $i : \Delta/\mathcal{X} \to \Delta$ est le
\uncertain{morphisme} d'\uncertain{incarnation} (\add{des topos} \uncertain{cat. associées}).

Si $\Delta/\mathcal{X} \xrightarrow{A^*} \mathcal{E}$ \struck{\ill{}} est un \og cofaisceau \fg{}
sur $\mathcal{X}$ (i.e. syst. à coeff. \emph{covariant} par rapport au simplexe variable de
$\mathcal{X}$), on lui associe $C^*(\mathcal{X}, A^*) = \ill{} : (\Delta \ill{}) \to \mathcal{E}$.

On passe de l'un à l'autre en remplaçant $\mathcal{E}$ par $\mathcal{E}^{\circ}$ ; c'est
compatible avec composition avec foncteurs $\mathcal{E} \to \mathcal{E}'$ qui commutent aux
sommes resp. produits.

\page{63}
\note{la page commence en cours d'argument : ce qui précède est avant le lot, la p. 62 étant étrangère.}
\marginal{(\uncertain{$k$ anneau commutatif} \ill{})}
Mais si $F$ \add{sur $X_*$} est loc. constant \add{(un $k$-Module)}, alors on montre
facilement que
\[
  C^*(X_*, F) \simeq \Gamma\bigl(X_*, \underbrace{(C^*_*|X_*)}_{C^*_{X_*}} \otimes_{k_*} F\bigr)
\]
\marginal{NB on \ill{} $C^*_* = C^*_{*k}$ \ill{} notations}
et on vérifie tout de suite que $C^*_{X_*} \otimes_{k_*} F$ est une résolution de $F$,
laquelle est d'ailleurs \emph{flasque} \marginal{à vérifier en passant} (que $F$ soit loc.
constant ou non, d'ailleurs). Donc, sans hypothèse de locale constance sur $F$,
$C^*(X_*, F)$ permet de calculer la coh. de $X_*$ à coeff. dans $F$ (i.e. on a un iso.
canon.
\[
  H^*\bigl(C^*(X_*, F)\bigr) \simeq H^*(X_*, F),
  \qquad
  C^*(X_*, F) \overset{\text{déf}}{=} \Gamma(X_*, C^*_{X_*} \otimes_k F),
\]
pour \struck{compatible} fonctoriel en $F$ (et \ill{} $(X_*, F)$ \ldots) et compatible
avec les cobords. De même, pour le calcul de l'\emph{hypercohomologie} de $X_*$ à coeff.
dans un complexe de $k$-Modules $F^*$. En termes de catégories dérivées, on a un
isomorphisme de foncteurs triangulés
\[
  R\Gamma(X_*, F^*) \simeq \underbrace{C^*(X_*, F^*)}_{\text{par diagonale associé aux deux degrés}}
\]
(venant de $C^*$ et de $F^*$).

\page{64}
Soit maintenant $A^*_*$ un \struck{$k$-module} complexe de cochaînes de $k$-modules
$\frac{1}{2}$simpliciaux
\[
  0 \to A^0_* \to A^1_* \to A^2_* \to \cdots
\]
et proposons-nous de définir les homomorphismes de ce complexe dans les
\[
  0 \to C^0_{*M} \to C^1_{*M} \to C^2_{*M} \to \cdots
\]
($M$ un $k$-module donné). Tout d'abord, les termes de degré $0$ \ill{} \ill{}
$A_* \to \ill{}$ \uncertain{(cf. (A))} sont les syst. d'hom.
\[
  A^n_* \xrightarrow{u_n} C^n_{*M},
\]
i.e. les syst. d'hom. de $k$-modules
\[
  A^n_n \xrightarrow{k_n} M .
\]
\struck{\ill{}} L'hom. composé
\[
  d_C \circ u_n : A^n_* \longrightarrow C^n_{*M} \longrightarrow C^{n+1}_{*M}
\]
correspond à l'hom. composé
\[
  A^n_{n+1} \xrightarrow{\;k_n \partial_{n+1}\;} M,
  \qquad
  A^n_{n+1} \xrightarrow{\partial_{n+1}} A^n_n \xrightarrow{k_n} M,
\]
tandis que l'hom. composé
\[
  u_{n+1} \circ d_A : A^n_* \longrightarrow A^{n+1}_* \longrightarrow C^{n+1}_{*M}
\]
correspond à l'hom. composé
\[
  A^n_{n+1} \longrightarrow M,
  \qquad
  A^n_{n+1} \xrightarrow{d^n_A} A^{n+1}_{n+1} \xrightarrow{k_{n+1}} M,
\]

\page{65}
et la condition que les $u_n$ définissent un hom. de complexes est la \og condition de
Stokes \fg{}
\[
  k_n \partial_{n+1} = k_{n+1} d :
\]

\begin{tikzcd}
A^n_{n+1} \arrow[r, "\partial_{n+1}"] \arrow[d, "d^n"'] & A^n_n \arrow[d, "k_n"] \\
A^{n+1}_{n+1} \arrow[r, "k_{n+1}"'] & M
\end{tikzcd}

Sous ces conditions, pour $X_*$ ($\in \mathrm{ob}\,\widehat{\Delta}$) \ill{} un hom.
\uncertain{déduit} de $u$
\[
  \Gamma(X_*, A^*_{X_*}) \Longrightarrow C^*(X_*, M_*) \simeq \Gamma(X_*, C^*_{*M}|X_*),
\]
qui représente un hom. dans la cat. dérivée
\[
  \Gamma(X_*, A^*_{X_*}) \longrightarrow R\Gamma_{X_*}(M)
\]
associé à $u_*$. On aura alors commutativité dans

\begin{tikzcd}
\Gamma(X_*, A^*_{X_*}) \arrow[r] \arrow[d] & {C^*(X_*, M) = \Gamma(X_*, C^*_M|X_*)} \arrow[d, "\wr"] \\
R\Gamma(X_*, A^*_{X_*}) \arrow[r] & R\Gamma(X_*, M)
\end{tikzcd}

Si $A^*$ est une résolution de $M_*$, et $k_0 : A^0_0 \to M$ induit l'identité sur
$M \subset A^0_0$, alors la dernière flèche horizontale est aussi un iso, et

\page{66}
on peut considérer la flèche
\[
  \underset{\textstyle \mathrm{Hom}(X_*, A^*)}{\Gamma(X_*, A^*_{X_*})}
  \longrightarrow C(X_*, M)
\]
déduite du syst. de Stokes--De Rham $(k_n)$ comme une réalisation de l'homom. canonique

\begin{tikzcd}
\Gamma(X_*, A^*_{X_*}) \arrow[r] & R\Gamma(X_*, A^*_{X_*}) \\
 & R\Gamma(X_*, M) \arrow[u, "\wr"']
\end{tikzcd}

\struck{On notera que sous les conditions de la prop. 1, il a}

\struck{Supposons}

\textbf{NB} L'hom.
\[
  \mathrm{Hom}(X_*, A^n_*) = \Gamma(X_*, A^n_{X_*}) \xrightarrow{u^n} C^n(X_*, M)
\]
associé à $k_n$ s'explicite ainsi : On considère, pour un \struck{\ill{}} élément
$f : X_* \to A^n_*$, la composante
\[
  f_n : X_n \to A^n_n
\]
et on prend
\[
  u^n(f) = k_n \circ f_n \in \mathrm{Hom}(X_n, M).
\]

Supposons maintenant \add{\uncertain{dans la $\frac{1}{2}$simpliciale}} que $A^*_*$
\struck{\ill{}} soit un complexe satisfaisant
\begin{enumerate}
  \item[a)] $A^i_n = 0$ si $i > n$ \quad (i.e. aussi $d A^n_n = Z^n_n$, i.e. $d(A^n_n) = 0$ si $n \geq 1$)
  \item[b)] $H^n(A^*_n) = 0$ si $n \geq 1$
  \item[c)] $H^{n-1}(A^*_n) = 0$ si $n \geq 2$.
\end{enumerate}
Alors on sait qu'un syst. $k_n : A^n_n \to M$ satisfaisant la condition de Stokes

\page{67}
est déterminé par la connaissance de
\[
  k_0 : A^0_0 = A_0 \longrightarrow M
\]
et $k_0$ correspond à un $(k_n)$ ssi
\[
  k_0\bigl(\partial_1(\mathrm{Ker}(A^0_1 \xrightarrow{d} A^1_1))\bigr) = 0
\]
(donc il y a un syst. d'intégrales de Stokes--De Rham \og universel \fg{} à valeurs dans
\[
  A_0 / \partial_1\bigl(\mathrm{Ker}(A^0_1 \xrightarrow{d} A^1_1)\bigr) ;
\]
le cas le plus intéressant est celui où $\partial_1(\mathrm{Ker}(A^0_1 \xrightarrow{d} A^1_1)) = 0$,
donc où le choix universel est fourni par $M = A_0$, ce qui sera le cas en particulier si
\begin{enumerate}
  \item[c')] $H^0(A^*_1) \xleftarrow{\;\sim\;} \varepsilon^*(A^0_0)$ \quad (où $\varepsilon : \Delta_1 \to \Delta_0$ \ldots)
\end{enumerate}
(ce qui est une variante de c) pour $n = 1$).
\marginal{ou $H^0(A^*_1) = \varepsilon^*(Z^0_0)$ \ill{} et \ill{} \ill{}}

Sous ces conditions, on sait donc qu'il existe un \emph{unique} hom. de complexes de
cochaînes $\frac{1}{2}$simpliciaux
\[
  A^*_* \longrightarrow C^*_{*M}
\]
qui sur $A^0_0 = A_0$ coïncide avec
\[
  k_0 : A^0_0 \longrightarrow C^0_{0M} = M
\]
(donc l'identité, i.e. toutes \uncertain{fonctions} \uncertain{pratiques}). On peut noter
que \struck{fait} cet hom. est à valeurs dans $C^{*!}_{*M}$ (les cochaînes non
dégénérées), i.e. \struck{fait}
\[
  A^*_* \longrightarrow C^{*!}_{*M} \simeq \Lambda^{*+1} \Phi_{*k} \otimes_k M,
\]
en se rappelant que
\[
  C^{n!}_{*M} = \bigcap_{\substack{\sigma : \Delta_m \to \Delta_n \\ m < n}}
  \mathrm{Ker}\bigl(C^n_{*M} \xrightarrow{\sigma^*} C^m_{*M}\bigr).
\]

\page{68}
d'où
\[
  \mathrm{Hom}(A^n_*, C^{n!}_{*M}) = \bigl\{ A^n_* \xrightarrow{f} C^n_{*M} \bigm|
  \sigma^* f = 0 \ \ \forall \sigma : \Delta_m \to \Delta_n,\ m < n \bigr\}.
\]

\begin{tikzcd}
A^n_* \arrow[r, "f"] \arrow[d, "\sigma^*_A"'] & C^n_* \arrow[d, "\sigma^*_C"] \\
A^m_* \arrow[r] & C^m_*
\end{tikzcd}

Or $A^n_* \to C^n_*$ est défini par
\[
  A^n_n \xrightarrow{k_n} M,
\]
son composé $A^n_* \to C^n_* \to C^m_*$ est défini par la
\[
  A^n_m \longrightarrow M
\]
qui n'est autre que le composé
\[
  A^n_m \xrightarrow{\sigma^A} A^n_n \xrightarrow{k_n} M,
\]
et \struck{la condition $f(A^n_*)$} \ill{} donc
\[
  f^n(A^n_*) \subset C^{n!}_* \iff k_n \circ \sigma^* = 0 \quad
  \forall \sigma : \Delta_n \to \Delta_m,\ m < n
\]
(il suffit \ill{} de le vérifier pour les \ill{} des dégénérescences) ; or cette
condition est toujours satisfaite si $A^i_m = 0$ si $i > m$, ce qu'on a supposé dans a).

\emph{Ex.} Soit $S$ un anneau \uncertain{à p.d.} $\ill{}$, \uncertain{v.d.}, $t \in J$,
considérons alors l'anneau différentiel \uncertain{simplicial à p.d.}
\[
  A^*_{*S} = \bigl(\Gamma^*(\Phi_{*S}) \otimes_S \Lambda^* \Psi_{*S}\bigr) / (T - t)_{pd},
\]
c'est une \emph{résolution} de $S$, et $A^i_j = 0$ si $i > j$, donc on a un syst.
d'intégrales de Stokes--De Rham à valeurs dans $S$, donc
\[
  A^*_* \xrightarrow{\;\sim\;} \bigl(\Gamma^*(\Phi_{*S}) \otimes_S \Lambda \Psi_{*S} / (T - t)_{pd}\bigr)
  \longrightarrow \Lambda^{*+1} \Phi_{*S}
\]
dont il faudrait étudier les propriétés multiplicatives \ldots{} \struck{\ill{}}

\page{69}
\begin{enumerate}
  \item[(1)] $k_X \xrightarrow{\varepsilon_X} A^*_X$, \quad $K_X = (K_n)_{n \geq 0}$,
  \quad $K_n : A_n \to k$, satisf.
  $\left\{\begin{array}{l} \text{Stokes} \\ K_0 \varepsilon_0 = \mathrm{id} \end{array}\right.$
  \item[(2)] Soit \struck{\ill{}} \uncertain{espace topologique} $X$, \add{\ill{}}
  faisceau différentiel $\varepsilon \to \Omega^*$ qui est une \emph{résolution} de $k_X$
  (p. ex. complexe de DR sur $X$ variété $C^{\infty}$, $k = \mathbb{R} = S$).
  \item[(3)] Un sous-ensemble $S(X)'$ du complexe singulier de $X$, tel que
  $S(X)' \to S(X)$ \struck{\ill{}} induise un isom. pour les groupes d'homologie
  \add{singulière} (p. ex. $X$ variété $C^{\infty}$, $S(X)'$ ens. des simplexes singuliers
  $C^{\infty}$) \add{(d'une base) et \ill{} \ill{} $S(X)$ \ill{} de $S$-complexes} \ill{}
  \struck{\ill{}} \ill{} $h$-augm.
  \item[(4)] $\forall \sigma \in S(X)'$, et \ill{}
  \[
    \sigma^* : \varinjlim_{U \supset \operatorname{Supp}\sigma} \Omega^*(U) \longrightarrow A^*_n
  \]
\end{enumerate}
\marginal{donc \ill{} le faisceau \ill{} \ill{} $S'_{X,k}$ \ill{} $S'^{*}_{X,k}$ \ill{}
$S'_{X,k}$ \ill{} une rés. \ill{}}
\note{« $\operatorname{Supp}\sigma$ » : lecture conjecturale d'une abréviation (« s--W »)
qui revient dans toute la page, sous les limites inductives.}

On suppose que si \ldots{} on a commutativité

\begin{tikzcd}
\varinjlim_{U \supset \operatorname{Supp}\sigma} \Omega^*(U) \arrow[r] \arrow[d, "\text{restriction}"'] & A^*_m \arrow[d] \\
\varinjlim_{V \supset \operatorname{Supp}\alpha^*(\sigma)} \Omega^*(V) \arrow[r] & A^*_n
\end{tikzcd}

\marginal{$\Delta_n \xrightarrow{\alpha} \Delta_m$, \quad $\sigma : \Delta_m \to S(X)' = \mathcal{X}$,
\quad $\alpha^*(\sigma) = \sigma \circ \alpha$}
\textbf{NB} $\operatorname{Supp}\alpha^*(\sigma) \subset \operatorname{Supp}\sigma$.

[\textbf{NB} On peut dire que les
\[
  \sigma \longmapsto \varinjlim_{U \supset \operatorname{Supp}\sigma} \Omega^*(U)
\]
forment un \struck{\ill{}} objet semi-simplicial $\Omega^*_{S(X)'}$ (en fait, un complexe
de $S$-modules $k$-augmentés) sur $S(X)' = \mathcal{X}$, et on se donne un hom. de
celui-ci dans le complexe augmenté $A^*_{*\mathcal{X}}$ sur $\mathcal{X} = S(X)'$ défini
par $A^*_{*}$.

\textbf{NB} \uncertain{Jusqu'ici} (1) $(A^*_*, \varepsilon_*, K_*)$ et (4) ne servent que
pour pouvoir définir un $\Omega^* \to S_X^{*\prime}$ \uncertain{compatible aux augm.}

\page{71}
\note{la première ligne de diagrammes est barrée de traits obliques ; on en donne ce qui se lit.}

\begin{tikzcd}[column sep=small, row sep=small, nodes={font=\scriptsize}]
R^*\Gamma(X, k) \arrow[r, "\varepsilon_{S^*_X}", "\sim"'] \arrow[d, "\wr"] & R^*\Gamma(X, S^*_X) \arrow[d, "\wr"] & \Gamma(X, S^*_X) \arrow[l, "\sim"'] \arrow[d, "\wr"] \\
R^*\Gamma(X, \Omega^*) \arrow[r, "\sim"] & R^*\Gamma(X, S'^*_X) & \Gamma(X, S'^*_X) \arrow[l, "\sim"'] \\
C^*(X, \Omega^*) \arrow[u] & &
\end{tikzcd}

\struck{avec $\Gamma(X, S^*_X) = C^*(S(X), k)$, $\Gamma(X, S'^*_X) = C^*(S(X)', k)$ ($S'^*_X$ \ill{} \ill{} \uncertain{homotopiquement} \ill{})}

\textbf{NB} \uncertain{la donnée} (4) \uncertain{jointe aux} $K_*$ définit un hom.
$\Omega^* \to S'^*_X$ \marginal{\uncertain{par $K$}} compatible avec les augmentations,
d'où carré commutatif

\begin{tikzcd}
\Omega^* \arrow[r] & S'^*_X \\
k \arrow[u] \arrow[r] & S^*_X \arrow[u]
\end{tikzcd}

\struck{(i) \ill{}} (d'isom. dans $D(X, S)$) d'où \uncertain{là} : un cube commutatif

\begin{tikzcd}[column sep=small, row sep=small, nodes={font=\scriptsize}]
C^*(X, \Omega^*) \arrow[r] \arrow[dd, "\wr"'] & {C^*(X, S'^*_X) \overset{\text{déf}}{=} C(S'_X, k)} \arrow[dd, "\wr"] & \\
 {[C(X, k)]} \arrow[rr, dashed] \arrow[u] & & {C^*(X, S^*_X) = C(S_X, k)} \arrow[ul, "\text{hypoth. } S'_X"'] \arrow[dd, "\wr"] \\
R^*\Gamma(X, \Omega^*) \arrow[r, "\sim"] & R^*\Gamma(X, S'^*_X) & \\
R^*\Gamma(X, k) \arrow[u, "\wr"] \arrow[rr, "\sim"] & & R^*\Gamma(X, S^*_X) \arrow[ul, "\wr"']
\end{tikzcd}

\note{la disposition du cube est redessinée de façon approximative : la flèche partant de
$[C(X, k)]$ est en pointillé et barrée en partie ; la flèche verticale de droite porte
« $\wr$ (Cartan, $S^*_X$ étant \og homotopiquement fin \fg) ».}

La base du cube est formée d'isom., car transformée par $R^*\Gamma$ d'un carré d'isos
dans une cat. dérivée ; la face \struck{\ill{}} du \uncertain{haut} est formée d'isos, car
\ill{} le sont déjà pour \uncertain{deux autres}, donc aussi pour la quatrième qui reste.

\page{73}
On travaille dans le topos $\mathcal{S}s = \widehat{\Delta}$, avec $k$ anneau comm.
\begin{enumerate}
  \item[] $\Phi_* = \Phi_{*k} = i_{0*}(k) = \Phi_{*\mathbb{Z}} \otimes_{\mathbb{Z}} k
  = \Phi_{*\mathbb{Z}} \otimes_{\mathbb{Z}_*} k_*$
\end{enumerate}
\struck{$T$} On le considère comme faisceau de $k$-modules, (pas comme faisceau
d'anneaux) muni d'une section $T$
\[
  \left\{\begin{array}{l}
    T \in \Gamma(\mathcal{S}s, \Phi_{*k}) \ \ [= \Gamma(\mathcal{S}s, C^{1,0}_*)] \\
    k_* \hookrightarrow \Phi_{*k}
  \end{array}\right.
\]
définissant une \emph{injection}.

On a l'\struck{algèbre} bigraduée \struck{\ill{}} différentielle à p.d. (diff. de bidegré
$-1, +1$ \ldots) \struck{(alg.} alg. alternée pour degré tot.,
\[
  C^{**}_{*k} = \mathcal{C}^{**}_{*k} = \Gamma^*_k \Phi_* \otimes_k \Lambda^* \Phi_*
\]
\marginal{Notations $\overline{\Phi}_*$ pour $\Lambda^1 \Phi_*$ et $\overline{x}$ pour
\og $1 \cdot x$ \fg{} \ill{} \ill{} \ill{}}
qui est une résolution de $k$. De façon précise,
\[
  C^{00}_{*k} = k_*,
\]
et en degré total $n > 0$, on trouve une suite exacte
\[
  0 \to \Gamma^n \Phi_* \to \Gamma^{n-1} \Phi_* \otimes \Lambda^1 \Phi_* \to
  \Gamma^{n-2} \Phi_* \otimes \Lambda^2 \Phi_* \to
\]
\[
  \cdots \to \Phi_* \otimes \Lambda^{n-1} \Phi_* \to \Lambda^n \Phi_* \to 0 \to 0 \cdots
\]
dont les termes \struck{sont} sont flasques.

On considère
\[
  dT \in \Gamma(\mathcal{S}s, C^{01}_*) \simeq \Gamma(\mathcal{S}s, \overline{\Phi}_*),
  \qquad dT = \overline{T},
\]
et \struck{l'homomorphisme} \ill{}
\[
  \overline{T}_\wedge : C^{**}_* \longrightarrow C^{**}_k ,
\]
de bidegré $(0, 1)$, \marginal{de carré nul, $\mathrm{Im}(\overline{T}_\wedge) = \mathrm{Ker}(\overline{T}_\wedge)$}
\add{anti}commutant à
\[
  d \quad (d(\overline{T} \cdot x) = \underbrace{d\overline{T}}_{=0} \cdot x - \overline{T}\, dx = -\overline{T}\, dx),
\]
car injectif \ill{} $C^{*0}$. On pose
\begin{align*}
  DR^{pq}_k = DR^{pq}_* &= \mathrm{Ker}\bigl(C^{p,q+1}_* \xrightarrow{\overline{T}_\wedge} C^{p,q+2}_*\bigr) \\
  &= \mathrm{Ker}\bigl(\Gamma^p \Phi_* \otimes \Lambda^{q+1} \Phi_* \to \Gamma^p \Phi_* \otimes \Lambda^{q+2} \Phi_*\bigr) \\
  &\simeq \text{\struck{$\mathrm{Ker}$}}\ \Gamma^p \Phi_* \otimes
  \underbrace{\mathrm{Ker}\bigl(\Lambda^{q+1} \Phi_* \to \Lambda^{q+2} \Phi_*\bigr)}_{\Lambda^q \Psi_*} ,
  \qquad \Psi_* = \Phi_* / k_* \\
  &\simeq \Gamma^p \Phi_* \otimes \Lambda^q \Psi_*
\end{align*}
\[
  DR^{**}_* = (DR^{pq}_*)_{(p,q) \in \mathbb{Z} \times \mathbb{Z}},
\]
avec la bigraduation et la différentielle (\uncertain{au signe près}) induites par
$C^{**}_*$ \add{(mais avec décalage \uncertain{de degré total})} \struck{\ill{}}, et une
structure d'\uncertain{anneau} \struck{pour le degré tot.} multiplicative \struck{qui est}
à p.d. qui \uncertain{sont} \uncertain{canoniques}

\page{74}
par la propriété que l'hom. épi
\[
  u = \text{\struck{$\overline{T}_\wedge$}} : C^{**}_* \xrightarrow{\;u\;} DR^{**}_*
  \overset{i}{\hookrightarrow} C^{**}_*
\]
soit compatible \struck{aux} structures (en fait, \emph{tous} ces structures). On aura aussi
\[
  \left|\begin{array}{l}
    i(x\, u(y)) = i(x)\, y \\
    i(u(y)\, x) = y\, i(x)
  \end{array}\right.
  \qquad i(x)\, i(y) = 0
\]
\struck{$i(u x\, u(y)) =$}

Notons que \uncertain{la} \uncertain{suite} exacte
\[
  0 \to DR^{p,q}_* \overset{i}{\longrightarrow} C^{p,q+1}_*
  \xrightarrow{\;(\overline{T}_\wedge) = u\;} DR^{p,q+1}_* \to 0,
  \qquad DR^{p,q+1}_* \subset C^{p,q+2}_*
\]
et que le noyau $DR^{p,q}_*$ \uncertain{est} flasque si $p \geq 0$ --- \uncertain{donc} \ill{}
si $X_* \in \mathrm{ob}\,\mathcal{S}s$, $\mathrm{Hom}(X_*, -)$ donne une suite exacte
(\uncertain{puisque} \struck{\ill{}} $C^{p,i}(X_*) \xrightarrow{\overline{T}} DR^{p,i}(X_*)$
est surjectif).

Posons donc $\Gamma = \mathrm{Hom}(X_*, -)$.
\[
  C^{**} = C^{**}_*(X_*) = C^{**}(X_*, k)
\]
est une $k$-algèbre bigraduée différentielle à p.d., avec \uncertain{élément}
\uncertain{marqué} $T \in C^{1,0}$,
\[
  \text{(A)} \quad
  \left\{\begin{array}{l}
    C^{00} \simeq k^{\pi_0(X_*)} \quad (= k \text{ si } X_* \text{ connexe}),
    \qquad = H^{00} \\
    H^{pq}(C^{**}) = 0 \ \text{ si } pq \neq 0,
  \end{array}\right.
\]
Posant $\overline{T} = dT \in C^{0,1}$, \ldots{} de plus
\[
  \text{(B)} \quad
  \left\{\begin{array}{l}
    0 \to C^{p,0} \xrightarrow{\overline{T}_\wedge} C^{p,1} \xrightarrow{\overline{T}_\wedge} C^{p,2} \to \cdots
    \to C^{p,q} \xrightarrow{\overline{T}_\wedge} C^{p,q+1} \to \\
    \text{est une suite exacte si } p > 0
  \end{array}\right.
\]
tandis que, pour $p = 0$, comme
\[
  \underset{\textstyle k}{\overset{0}{\Lambda} \Phi_*}
  \xrightarrow{\overline{T}_\wedge} \overset{1}{\Lambda} \Phi_*
  \xrightarrow{\overline{T}_\wedge} \overset{2}{\Lambda} \Phi_* \to
\]
est une suite exacte fournissant une résolution flasque de $k_*$, on trouve
\[
  H^n\bigl(0 \to C^{0,0} \xrightarrow{\overline{T}_\wedge} C^{0,1} \to \cdots \to
  C^{0,q} \xrightarrow{\overline{T}_\wedge} C^{0,q+1} \to \cdots\bigr)
  \simeq \left\{\begin{array}{ll}
    0 & \text{si } n \leq 1 \\
    H^{n-1}(X_*, k) & \text{si } n \geq 2
  \end{array}\right.
\]
\note{$H^n$ pour $n \le 1$ nul et $H^{n-1}(X_*,k)$ pour $n \ge 2$ : tel quel sur la page.}

\page{75}
Posant
\[
  DR^{pq} = DR^{pq}(X_*) = \mathrm{Ker}\bigl(C^{p,q+1} \xrightarrow{\overline{T}_\wedge} C^{p,q+2}\bigr),
  \qquad
  DR^{**} = (DR^{pq})_{p,q \in \mathbb{Z}},
\]
on aura donc
\[
  C^{p,q} \xrightarrow{\;\overline{T}_\wedge =: u\;} DR^{pq}, \ \subset C^{p,q+1},
  \quad \text{épi si } p \neq 0,
  \qquad \text{d'où} \quad
  C^{**} \xrightarrow{\;u\;} DR^{**} \overset{i}{\hookrightarrow} C^{**}
\]
\marginal{$\mathrm{Im}\, i = \mathrm{Ker}\, u$, \quad $i \circ u = \overline{T}_\wedge(\cdot)$}
\[
  \text{(C)} \quad i(x) \cdot i(y) = 0
\]
(il en résulte que $i(u(\ldots))$ \ldots{} $i(u\, DR) = 0$ donc $i(DR^{p,q}) \times i(DR^{p',q'}) = 0$
si $p \neq 0$, $p' \neq 0$)
\marginal{donc (C) implique $i(DR^{0,p}) \cdot i(DR^{0,q}) \neq 0$ \ill{}}
et on aura sur $DR^{**}$ une structure d'algèbre bigraduée différentielle à p.d.
\add{commutative, avec \uncertain{exposants} \uncertain{usuels} (\uncertain{alternée} etc.)}
\struck{telle que}
\[
  u : C^{**} \longrightarrow DR^{**}
\]
soit compatible \ldots{} \uncertain{toutes} \uncertain{ces} structures
\add{et $\left\{\begin{array}{l} i(x\, u(y)) = i(x)\, y \\ i(u(y)\, x) = y\, i(x) \end{array}\right.$}
\uncertain{Cette} structure sur $DR^{**}$ est presque déterminée par ces propriétés, plus
celle que l'\add{anti}commutant $DR^{**} \to C^{**}$ est de bidegré $0, 1$ et commutant à
la différentielle --- à la détermination près des produits $xy$ \add{\ill{}} et des
puissances divisées $x^{(i)}$ ($i \geq 2$) \uncertain{prenant} $x, y$ de degré
complémentaire, (auquel cas on ne \uncertain{peut} \ill{} pas \uncertain{nécessairement}
\uncertain{les} écrire sous la forme $u(x')$ \add{$u(y')$} \uncertain{tels} dans $C^{**}$).
\add{$(\mathrm{Im}\, u) \times DR$} \struck{\ill{}} \add{et $DR \times \mathrm{Im}\, u \to$}
\struck{\ill{}}

(\textbf{NB} Il est évident que \uncertain{la} \struck{\ill{}} \add{\uncertain{définition}
\uncertain{sur} $DR \times DR$} \add{à \uncertain{p.d.} \ill{}} multiplicat., \ill{} satisfait
\[
  u(x \cdot y) = u(x)\, u(y), \qquad i(x\, u(y)) = i(x)\, y, \qquad i(u(y)\, x) = y\, i(x)
\]
(\struck{\ill{}} la première relation \uncertain{équivaut}, en \uncertain{appliquant} $i$,
à $\overline{T} \cdot (xy) = i(u x \cdot u y)$, i.e. en \uncertain{multipliant} à
\uncertain{droite}
\[
  \overline{T} \cdot (xy) = (\overline{T} x)\, y
\]
donc il suffit de vérifier l'associativité et l'unicité de
\[
  \mathrm{Im}\, u \times DR \to DR \quad \text{et} \quad DR \times \mathrm{Im}\, u \to DR
\]
satisfaisant les rel. 2 et 3. Or pour la $2^{\text{e}}$, il suffit de
\uncertain{montrer} que
\begin{enumerate}
  \item[a)] $i(x)\, y \in$ \struck{$\mathrm{Ker}$} $DR = \mathrm{Ker}(\overline{T}_\wedge)$
  \add{$= xy$}, car $\overline{T}(xy) = (\overline{T} x)\, y = 0$ car $\overline{T} x = 0$
  \uncertain{puisque} $x \in DR$
  \item[b)] $i(x)\, y$ \add{$= xy$} ne dépend que de $x \in DR$ \uncertain{et} de
  $u y \in DR$, car si $u(y) = 0$ i.e. $\overline{T} y = 0$, i.e. $y \in DR$, \ldots{}
  $xy = 0$ --- Idem pour 3.
\end{enumerate}

\page{76}
\textbf{D} \ldots, \uncertain{il} \uncertain{suffit} \uncertain{que} $u(C^{**+}) \subset
DR^{**+}$ \uncertain{des} \uncertain{op.} \og puiss. div. \fg{} \add{\uncertain{marquées}}
\uncertain{données} par la condition
\[
  u(x)^{(i)} = u(x^{(i)})
\]
car il suffit de prouver que si $x'$ est tel que $u x = u x'$, i.e. $x' = x + y$,
$\overline{T} \cdot y = 0$, alors $u(x'^{(i)}) = u(x^{(i)})$, i.e.
$\overline{T}_\wedge x'^{(i)} = \overline{T} x^{(i)}$
\[
  x'^{(i)} = x^{(i)} + x^{(i-1)} y + x^{(i-2)} y^{(2)} + \cdots + y^{(i)},
\]
il suffit de prouver que $\overline{T}\, y^{(i)} = 0$ pour $i \geq 1$, ce qui résulte du
fait \struck{\ill{}}
\begin{enumerate}
  \item[(D)] $\underset{\textstyle \mathrm{Ker}(\overline{T}_\wedge(\cdot))}{DR^{**+}} \subset C^{**}$
  \struck{est stable par} \add{\uncertain{est}} les puiss. div. (dans $C^{**}$)
  \add{$y^{(i)}$ ($i \geq 2$) $\Rightarrow$ \uncertain{nulles}}
\end{enumerate}
\struck{i.e. si $y \in C^{**+}$, $\overline{T} \cdot y = 0 \Rightarrow \overline{T} \cdot y^{(i)} = 0$
$\forall i \geq 1$,}

(ceci \uncertain{se} \uncertain{démontre} résulte de la \uncertain{construction} pour
$DR^{**}_* \subset C^{**}_*$, où loc. $y$ est de la forme $\overline{T} z$, donc
$\overline{T}\, y^{(i)} = \overline{T}^{i} z^{(i)}$, \uncertain{qui} est donc nul si
$i \geq 2$.)

\struck{(\ill{}) $i(\ldots\overline{T}\ldots DR^{\ldots} \subset C\ldots) \ldots$}

Ainsi, les seules données supplémentaires sur $C^{**}$ pour \add{déterminer}
\struck{définir} les structures habituelles sur $DR^{**}$ sont
\[
  \text{(E)} \quad
  \left\{\begin{array}{lll}
    \text{multiplication} & DR^{0,p+} \times DR^{0,q+} \longrightarrow DR^{0,p+q+} & (p, q \geq 1) \\
    \text{puiss. division} & DR^{0,p} \xrightarrow{\;\gamma^i\;} DR^{0,pi} & (i \geq 2)
  \end{array}\right.
\]
\marginal{(\uncertain{marquées})}
\note{dans (E), « DR » est chaque fois écrit sur un mot biffé illisible.}

\struck{(\textbf{NB} \ill{} $Z^{0,1} \times Z^{0,q+1} \to \ldots$ pour $p$ ou $q = 1$,
$Z^{0,1}$ est engendré par \ldots{} $Z^{0,p+1}$)}

\uncertain{plus} les axiomes \uncertain{auxquels} il faut \uncertain{pour} \uncertain{assurer}
que sur $DR^{**}$, \ldots{} les propriétés habituelles pour les opérations.

\page{77}
\textbf{Pb} Trouver des \uncertain{splittings} \emph{canoniques} des suites exactes
\[
  \begin{array}{ccc}
    0 & & \\
    \downarrow & & \\
    \Gamma^p \Phi_* \otimes \Lambda^q \Psi_* & & \cdots \to \Lambda^{q+1} \Phi_* \\
    \uparrow & \text{?} & \\
    \Gamma^p \Phi_* \otimes \Lambda^q \overline{\Phi}_* & & \\
    \uparrow & & \\
    \Gamma^p \Phi_* \otimes \Lambda^{q-1} \Psi_* & & \cdots \to \Lambda^q \Phi_* \\
    \uparrow & & \\
    0 & &
  \end{array}
\]
pour $p \geq 1$.
\note{la flèche courbe marquée « ? » va du terme du haut vers celui du milieu : c'est le
scindage cherché.}

\textbf{NB} Comme $\Gamma^p \Phi_* \otimes \Lambda^q \Psi_*$ est un $k$-Module projectif
\add{($p \geq 1$)}, il y a un splitting. On veut un choix simple. Notons que ce cas $q = 0$
est trivial (car $\Lambda^{q-1} \Psi_* = 0$).

\emph{Cas premier} $p = 1$, $q = 1$.

La suite exacte est
\[
  0 \to \underset{\textstyle \Phi_*}{\Phi_* \otimes k} \longrightarrow
  \Phi_* \otimes \overline{\Phi}_* \longrightarrow \Phi_* \otimes \Psi_* \to 0
\]
On note que
\[
  \mathrm{Hom}_{k\text{-Mod}}\bigl(\Phi_* \otimes \Psi_*, \underset{\textstyle i_{0*}(k)}{\Phi_*}\bigr)
  \simeq \mathrm{Hom}_{k\text{-Mod}}\bigl((\Phi_* \otimes \Psi_*)_0, k\bigr) = 0,
  \qquad \text{car } (\Phi_* \otimes \Psi_*)_0 = \Phi_0 \otimes \Psi_0 = 0,
\]
donc splitting \emph{unique}
\[
  \Phi_* \otimes \overline{\Phi}_* \simeq (\Phi_* \otimes \Psi_*) \oplus \Phi_*
\]
Soit \struck{\ill{}}
\[
  s : \Phi_* \otimes \Psi_* \longrightarrow \Phi_* \otimes \overline{\Phi}_*
\]
\note{la page s'arrête ici ; la définition de $s$ n'est pas donnée.}

\section{Faisceaux tests du topos De Rham}
\note{titre encadré de sa main, en haut à droite de la p. 79.}

\page{79}
1. Soit $X$ un topos, et
\[
  \mathbb{Z} \xrightarrow{\;u\;} \Phi
\]
une immersion de $\mathbb{Z}$ dans un faisceau $\Phi$, tel que
\begin{enumerate}
  \item[a)] $\mathrm{Coker}\,\Phi$ \note{sic, pour le conoyau de $u$} est $\mathbb{Z}$-plat
  \item[b)] Pour tout $i > 0$, $\Lambda^i_{\mathbb{Z}} \Phi$ et $\Gamma^i_{\mathbb{Z}} \Phi$
  sont flasques, et le restent par tensorisation avec tout faisceau abélien.
  \struck{\ill{}}
\end{enumerate}
En particulier, on en conclut \struck{\ill{}}
\begin{enumerate}
  \item[b')] Si $i, j \geq 0$ et $i + j > 0$, alors \struck{$\Lambda$}
  $\Gamma^i_{\mathbb{Z}} \Phi \otimes \Lambda^j_{\mathbb{Z}} \Phi$ est flasque.
\end{enumerate}
De b') on conclut \struck{\ill{}} \add{(par récurrence et dévissage)}, posant
\[
  \mathrm{Coker}\, i = \Psi
\]
\begin{enumerate}
  \item[$\alpha$)] \add{que} les $\Gamma^i_{\mathbb{Z}} \Phi \otimes \Lambda^j_{\mathbb{Z}} \Psi$
  sont flasques pour $i > 0$.
\end{enumerate}
De ceci on conclut \struck{\ill{}}
\begin{enumerate}
  \item[$\beta$)] Posant $K^{\cdot\cdot}_u = \Gamma^*_{\mathbb{Z}} \Phi \otimes \Lambda^*_{\mathbb{Z}} \Psi$,
  avec l'\uncertain{opér.} diff. \struck{\ill{}} et compatible aux puissances divisées, qui,
  sur \struck{\ill{}} $\Lambda \Psi$ est nulle et sur $\Phi$ est l'hom. canonique
  $\Phi \to \Psi$, et définissant
  \[
    \mathbb{Z}\{T\} \longrightarrow \mathcal{H}^{*0}(K^{**}_u)
  \]
  par $T \mapsto 1 \in \mathbb{Z} = \mathrm{Ker}(\Phi \xrightarrow{d} \Psi)$, et prolongeant
  par la condition d'être compatible aux p. div., on trouve que $K^{\cdot\cdot}_u$ est une
  \emph{résolution} de $\mathbb{Z}\{T\}$.
\end{enumerate}
Utilisant encore a), \add{et b)}, on trouve
\begin{enumerate}
  \item[$\gamma$)] Pour tout faisceau d'anneaux comm. $k$ sur $X$,
  $K^{\cdot\cdot}_{u,k} \simeq K^{\cdot\cdot}_u \otimes k \simeq \Gamma^*_k \Phi_k \otimes \Lambda^*_k \Psi_k$
  est une résolution de $k\{T\} \simeq \mathcal{H}^{\cdot 0}(K^{\cdot\cdot}_{u,k})$, et les
  $\Gamma^i_k \Phi_k \otimes \Lambda^j_k \Psi_k$ sont flasques pour $i > 0$.
\end{enumerate}

\page{80}
On va graduer $K^{\cdot\cdot}_{u,k}$ par a) le degré extérieur venant de
$\Lambda^*_k \Psi_k$ (deg ext $\Gamma^i \otimes \Lambda^j = j$) et b) le \emph{degré
total} $i + j$. Ainsi $d$ augmente le degré extérieur par $1$, et invarie le degré total
\struck{\ill{}} (le degré complémentaire étant deg comp $\Gamma^i \otimes \Lambda^j = i$,
il est diminué de $1$ par $d$\ldots) Muni des degrés total (ou complémentaire)
$K^{\cdot\cdot}_{u,k}$ est une Algèbre graduée sur l'Anneau gradué $k\{T\}$.

Conséquences (de ce qui précède de suite)

\emph{Théorème} Soit $K^{\cdot\cdot}(X, u, k) = \Gamma(X, K^{\cdot\cdot}_{u,k})$.
Alors on a un hom. can. (du fait que $K^{\cdot\cdot}_{u,k}$ \struck{\ill{}} est une
résolution de $k\{T\}$)
\[
  H^{**}(X, u, k) \overset{\text{déf}}{=} H^{**}\bigl(K^{\cdot\cdot}(X, u, k)\bigr)
  \longrightarrow H^{**}(X, k\{T\}) \simeq H^*(X, k) \otimes_{\mathbb{Z}} \mathbb{Z}\{T\}
\]
\note{le dernier membre est écrit $\mathbb{Z}\{T\}$ et non $k\{T\}$ ; tel quel.}
et ceci est un isomorphisme pour les bidegrés $(i, j)$ tels que \struck{\ill{}}
$j \leq i$, tandis que $H^{i,j}(X, u, k) = 0$ (et pas un iso.) pour $j > i$.

\drawing{dans la marge gauche, un quadrant d'axes avec des points disposés en escalier
sous la diagonale, figurant les bidegrés $(i,j)$, $j \le i$, où le morphisme est un
isomorphisme.}

\emph{Corollaire plus général.} Soit $f : X \to Y$ un morphisme de topos. Alors
\[
  f_* K^{\cdot\cdot}_{u,k} \longrightarrow Rf_*(k\{T\}) \simeq Rf_*(k) \otimes_{\mathbb{Z}} \mathbb{Z}\{T\}
\]
\note{la page s'arrête ici ; l'énoncé du corollaire se poursuit au-delà du lot.}

\end{document}
